\documentclass[12pt,a4paper]{article}
\usepackage[a4paper,top=15mm,bottom=15mm,left=16mm,right=16mm]{geometry}
\usepackage[T1]{fontenc}
\usepackage{amsmath}
\usepackage{amssymb}
\usepackage{multicol}

\pagestyle{empty}
\setlength{\parindent}{0pt}
\setlength{\parskip}{3pt}
\setlength{\columnsep}{9mm}

\newcommand{\topic}[1]{\par\addvspace{8pt}{\bfseries #1}\par\nobreak\vspace{1pt}}
\newcommand{\ans}[1]{\par\hangindent=1.8em\hangafter=1\noindent\makebox[1.8em][l]{\textbf{#1}}}

\begin{document}

\begin{center}
{\Large\bfseries Fractions, decimals and percentages}\\[2pt]
{\large Revision questions: answers}
\end{center}
\hrule\vspace{4pt}

\begin{multicols}{2}

\topic{Simplifying fractions}
\ans{1} $\frac{18}{24}=\frac34$
\ans{2} (a) $\frac34$\quad (b) $\frac23$\quad (c) $\frac7{10}$

\topic{Ordering fractions}
\ans{3} (a) $\frac7{12},\ \frac58,\ \frac34$\quad (b) $\frac{13}{20},\ \frac7{10},\ \frac34,\ \frac45$
\ans{4} Sam, because $\frac25=\frac{16}{40}$ is more than $\frac38=\frac{15}{40}$

\topic{Mixed numbers and improper fractions}
\ans{5} (a) $2\frac34$\quad (b) $\frac{11}4$
\ans{6} (a) $\frac{17}3$\quad (b) $5\frac16$

\topic{Adding and subtracting fractions}
\ans{7} (a) $\frac8{15}$\quad (b) $\frac{13}{28}$\quad (c) $\frac{29}{24}=1\frac5{24}$
\ans{8} $\frac7{20}$

\topic{One number as a percentage of another}
\ans{9} (a) 24\%\quad (b) 45\%\quad (c) 17.5\%
\ans{10} Priya (85\%, compared with Leo's 82\%)

\topic{Percentages without a calculator}
\ans{11} (a) 47\quad (b) 43\quad (c) 6.5\quad (d) 54
\ans{12} £36

\topic{Converting between fractions, decimals and percentages}
\ans{13} $\frac9{25}$, 0.36, 36\%
\ans{14} $\frac35=0.6=60\%$\\
$\frac9{20}=0.45=45\%$\\
$\frac1{25}=0.04=4\%$

\topic{Ordering fractions, decimals and percentages}
\ans{15} (a) $\frac38$, 0.4, 41\%\quad (b) 0.3, 33\%, $\frac13$
\ans{16} Shop C ($0.35=35\%$, $\frac13\approx33.3\%$)

\topic{Fractions of amounts}
\ans{17} 36
\ans{18} (a) 24\quad (b) 392
\ans{19} 480

\topic{Percentages of amounts}
\ans{20} (a) 49\quad (b) 8
\ans{21} (a) 19.55\quad (b) £108

\topic{Percentage change}
\ans{22} (a) 24\% increase\quad (b) 15\% decrease
\ans{23} 25\%

\topic{Simple interest}
\ans{24} £510, £520, £530, £540
\ans{25} 5 years

\topic{Calculating with mixed numbers and negatives}
\ans{26} (a) $3\frac{11}{12}$\quad (b) $2\frac7{10}$\quad (c) 4
\ans{27} (a) 3\quad (b) $-\frac35$\quad (c) 2

\topic{Word problems with mixed numbers}
\ans{28} $2\frac7{12}$\,m
\ans{29} 8 glasses
\ans{30} $21\frac23$\,m$^2$

\topic{Recurring decimals to fractions}
\ans{31} $10n=4.444\ldots$, $9n=4$, $n=\frac49$
\ans{32} (a) $\frac4{11}$\quad (b) $\frac56$

\topic{Percentage increase and decrease with multipliers}
\ans{33} (a) 1.08\quad (b) 0.7\quad (c) 1.025
\ans{34} (a) £86.40\quad (b) 1100

\topic{Finding the original amount}
\ans{35} £80
\ans{36} £70
\ans{37} £45

\topic{Problem solving}
\ans{38} Shop B, by 15p (Shop A £60, Shop B £59.85)
\ans{39} The 30\% was taken off the original price, not off £280. $280\div0.7=$ £400
\ans{40} For example $\frac{13}{20}=0.65$, $\frac{31}{50}=0.62$ and $\frac23=0.\dot6$, which are all between $\frac35=0.6$ and $\frac7{10}=0.7$

\end{multicols}

\end{document}
